\section{Introduction}
\label{sec:intro}

% PAGE TARGET: ~5 pages.
% Day 2 (5/26).
% NOTE: CUED template puts prior-work / background INSIDE the Introduction,
% not in a separate Background chapter. Tell the reader WHY we are doing
% this work.

\subsection{Motivation}
% - Codicology / paper history relies on laid-line density and chain-line
%   spacing to fingerprint papermaking moulds.
% - Manual counting under raking light is slow, subjective, and risks damage
%   if done at scale.
% - Multispectral imaging (MSI) is now standard in major libraries
%   (Cambridge UL, Bodleian, BnF), creating large image stacks that
%   could be analysed automatically -- but no robust pipeline exists.

\subsection{Problem Statement}
% - Given a calibrated image stack (DRP or single-band TX), recover:
%     (i)  laid-line density (lines/cm),
%     (ii) wire diameter (FWHM, mm),
%     (iii) per-region orientation,
%   with quantified uncertainty.

\subsection{Prior Work}
% Brief literature pass (one page).
% - Hiary \& Ng (2007), van der Lubbe et al., recent Bodleian / BnF efforts.
% - Common weaknesses:
%   (a) orientation must be supplied,
%   (b) sensitive to scan noise,
%   (c) no uncertainty quantification,
%   (d) no per-region (patchwise) handling for curved or warped folios.

\subsection{Contributions}
% Three contributions to claim:
% 1. A staged DRP-based pipeline (Stage 0--5) needing no per-folio tuning
%    beyond optical scale.
% 2. A synthetic phantom characterisation framework: period, SNR, angle,
%    wire-width sweeps.
% 3. A 9-folio real-world benchmark with manual GT cross-check on two
%    folios, achieving sub-1\% density error in the calibrated comparison.

\subsection{Report Structure}
% One paragraph mapping sections 2--6.
